% Generated from rules.md by scripts/render_statement.py.
% Public statement only; edit the Markdown source, then render again.
\documentclass[11pt,a4paper]{article}
\usepackage[margin=20mm]{geometry}
\usepackage{fontspec,xeCJK}
\setmainfont{Liberation Serif}
\setsansfont{DejaVu Sans}
\setmonofont[Scale=0.83]{DejaVu Sans Mono}
\setCJKmainfont[AutoFakeBold=2]{AR PL SungtiL GB}
\setCJKsansfont[AutoFakeBold=2]{AR PL SungtiL GB}
\setCJKmonofont{AR PL SungtiL GB}
\usepackage{amsmath,amssymb}
\usepackage{longtable,booktabs,array,calc}
\usepackage{graphicx,xcolor}
\usepackage{float}
\floatplacement{figure}{H}
\usepackage[font=small,labelsep=quad]{caption}
\renewcommand{\figurename}{图}
\newcounter{none} % Pandoc tables set \LTcaptype{none}; caption needs that counter.
\makeatletter
\@ifundefined{KV@Gin@alt}{\define@key{Gin}{alt}{}}{}
\newsavebox\pandoc@box
\newcommand*\pandocbounded[1]{\sbox\pandoc@box{#1}%
  \Gscale@div\@tempa{\textheight}{\dimexpr\ht\pandoc@box+\dp\pandoc@box\relax}%
  \Gscale@div\@tempb{\linewidth}{\wd\pandoc@box}%
  \ifdim\@tempb\p@<\@tempa\p@\let\@tempa\@tempb\fi
  \ifdim\@tempa\p@<\p@\scalebox{\@tempa}{\usebox\pandoc@box}\else\usebox{\pandoc@box}\fi}
\makeatother
\usepackage[unicode,hidelinks]{hyperref}
\usepackage{titlesec,enumitem}
\setcounter{secnumdepth}{0}
\titleformat{\section}{\large\bfseries}{}{0pt}{}
\titlespacing*{\section}{0pt}{11pt}{5pt}
\setlength{\parindent}{0pt}
\setlength{\parskip}{4pt}
\setlength{\emergencystretch}{2em}
\renewcommand{\arraystretch}{1.14}
\setlength{\LTpre}{4pt}
\setlength{\LTpost}{4pt}
\linespread{1.08}
\clubpenalty=10000
\widowpenalty=10000
\setlist{leftmargin=1.6em,itemsep=4pt,topsep=3pt}
\providecommand{\tightlist}{\setlength{\itemsep}{3pt}\setlength{\parskip}{0pt}}
\hypersetup{pdftitle={SOP-UQ 详细规则},pdfauthor={SOP-UQ},pdfsubject={SOP-UQ public document}}
\begin{document}
\begin{center}\fontsize{15}{20}\selectfont\bfseries SOP-UQ 详细规则\end{center}
\vspace{2pt}
本文件是题面（\texttt{docs/problem\_statement.md}）的补充，列出考试流程、提交、费用、数据与工具、评分和
baseline 的具体规则。符号含义与题面相同。

\section{1. 考试流程}\label{ux8003ux8bd5ux6d41ux7a0b}

\begin{enumerate}
\def\labelenumi{\arabic{enumi}.}
\tightlist
\item
  正式测试共 360 题。考官事先用固定 agent 生成每题的初始答案
  \(a_i\)、调用记录和共享采样，所有考生使用同一份。提交的初始答案必须与考官记录一致；考官记录为失败的题，不能标为
  \texttt{completed}。
\item
  其中 72 道是交互题，可以调用普通工具和 oracle；其余 288
  道不能调用任何工具，只评估查询前风险
  \(p_i\)。哪些题是交互题由服务告知。
\item
  交互题调用 oracle 之前，必须先调用 \texttt{checkpoint\_initial}
  锁定本题的 \(p_i\)（也可以锁定为失败或缺失状态）。只要 oracle
  还没有成功返回过结果，仍可以修改 \(p_i\)，但修改后需要重新锁定；oracle
  成功返回之后，\(p_i\)
  不能再改。如果正式考试有多个初始模型，同一题所有模型的 \(p_i\)
  一起锁定。
\item
  没有调用过 oracle 的题，\(p_i\)
  可以在后续提交版本中修改，以最后一个版本为准。
\item
  预测时可以利用其他题目的信息，可以在考试中更新方法参数，也可以更换后续求解的模型或算法。方法和配置的变化要记录下来，便于复现。
\item
  调用记录、费用和额度按''考生 +
  试卷''统计。更换方法、模型、运行或提交版本都不会重置。
\end{enumerate}

\section{2. 提交与交卷}\label{ux63d0ux4ea4ux4e0eux4ea4ux5377}

\begin{enumerate}
\def\labelenumi{\arabic{enumi}.}
\tightlist
\item
  每名考生最多提交 5
  个不同版本。每个版本包含全部题目；每道交互题的终局记录必须在该题最后一次调用之后提交。
\item
  内容相同且其间没有新调用时，重复上传不计次数；格式错误或与服务端记录不一致的包会被拒绝，也不计次数。
\item
  第 5
  个版本被接收后自动交卷；也可以在提交至少一个版本后提前交卷；到截止时间自动交卷，此前没有提交过版本的记为未提交。
\item
  计分时，调用过 oracle 的题使用锁定的 \(p_i\)；其余题的
  \(p_i\)，以及所有交互题的 \(b_i\)、\(q_i\)、\(r_i\)
  和决策，使用最后一个版本。费用按交卷时的累计值计算，包括最后一个版本之后发生的调用。
\item
  考试期间每次提交只返回''已接收''，不返回分数、错误率或任何与标签有关的信息。所有考生交卷或考试截止后，考官统一公布成绩，并永久关闭本场考试的报名和修改；公布过成绩的测试数据不能再用于下一批考生。
\item
  需要提交：可运行的代码、配置、训练得到的模型和校准文件；用所附客户端生成的
  \texttt{submission.zip}（包含 \texttt{manifest.json} 和
  \texttt{results.csv}）；technical
  report。初始答案和已锁定的风险必须与服务端记录一致，终局答案和风险必须基于该版本提交时的最新调用记录；费用和调用次数以服务端记录为准。公开报告与给考官的复现材料分开提交。
\end{enumerate}

\section{3. 费用与预算}\label{ux8d39ux7528ux4e0eux9884ux7b97}

{\small\def\LTcaptype{none} % do not increment counter
\begin{longtable}[]{@{}ll@{}}
\toprule\noalign{}
操作 & 每次费用 \\
\midrule\noalign{}
\endhead
\bottomrule\noalign{}
\endlastfoot
普通工具调用 & 0.01 \\
用 \texttt{independent} 方式读取可能有噪声的字段（第 5 节） & 0.03 \\
调用 oracle 查询一个字段 & 0.10 \\
转人工（DEFER） & 0.50 \\
\end{longtable}
}

\begin{enumerate}
\def\labelenumi{\arabic{enumi}.}
\tightlist
\item
  正式测试中，每道题的工具和 oracle 费用合计不超过
  2.00，其中普通工具最多 24 次，oracle 最多 16
  次；考官生成初始答案时已发生的费用也计入。
\item
  网络重试同一请求不重复收费，新的请求照常收费；对已经查过的 oracle
  字段再次查询，会免费返回原结果。
\item
  2.00
  的预算主要用于防止资源滥用，对现有的合理查询路径一般不构成约束，因此不能据此说明方法在预算紧张时的表现。研究预算约束时，可以使用
  0.50 或 1.00
  的公开配置，在相同配置下重新采集参考轨迹和校准数据，并单独报告。
\end{enumerate}

\section{4. dev 练习}\label{dev-ux7ec3ux4e60}

\begin{enumerate}
\def\labelenumi{\arabic{enumi}.}
\tightlist
\item
  \texttt{dev} 的 120 道题用于练习，其中 60 道为偏移条件。
\item
  练习时 oracle
  不限次数、不计入预算、不需要锁定风险；提交后立即返回成绩，不计入排名，也不占测试的提交次数和额度。名义费用仍会记录，便于比较策略。
\item
  练习中每道题的普通工具仍限 24 次，用完可以开始新的练习。
\item
  练习入口不能访问测试题。
\end{enumerate}

\section{5. 数据与工具}\label{ux6570ux636eux4e0eux5de5ux5177}

\textbf{有噪声的字段。}
案例的输入字段是准确的。工具返回的结果中，只有下表中的字段可能与真实值不同，其他返回值都准确。oracle
只能查询这些字段，每次返回一个字段在指定业务阶段的真实值；危险品的四个分数在
oracle 中分别叫 \texttt{sds}、\texttt{handling}、\texttt{transportation}
和 \texttt{disposal}。客服任务要先调用
\texttt{performTechnicalDiagnostics} 才能查询 \texttt{pre\_repair}
阶段的值，调用 \texttt{executeTroubleshooting} 之后才能查询
\texttt{post\_repair} 阶段的值。

{\small\def\LTcaptype{none} % do not increment counter
\begin{longtable}[]{@{}
  >{\raggedright\arraybackslash}p{(\linewidth - 2\tabcolsep) * \real{0.2000}}
  >{\raggedright\arraybackslash}p{(\linewidth - 2\tabcolsep) * \real{0.8000}}@{}}
\toprule\noalign{}
\begin{minipage}[b]{\linewidth}\raggedright
业务域
\end{minipage} & \begin{minipage}[b]{\linewidth}\raggedright
可能有噪声的字段（返回该字段的工具）
\end{minipage} \\
\midrule\noalign{}
\endhead
\bottomrule\noalign{}
\endlastfoot
危险品 &
\texttt{sds\_label\_score}、\texttt{handling\_score}、\texttt{transportation\_score}、\texttt{disposal\_score}（四个
\texttt{calculate\_*\_score} 工具） \\
飞机部件检查 &
\texttt{aircraft\_ready}（\texttt{VerifyAircraftClearance}）、\texttt{mechanical\_inspection\_result}（\texttt{VerifyMechanicalComponents}）、\texttt{electrical\_inspection\_result}（\texttt{VerifyElectricalSystems}）、\texttt{cross\_check\_response}（\texttt{CrossCheckSpecifications}） \\
客服排障 &
\texttt{latency}、\texttt{jitter}、\texttt{bandwidth}（\texttt{performTechnicalDiagnostics}
和 \texttt{executeTroubleshooting} 返回的服务指标） \\
\end{longtable}
}

\textbf{观测条件。}
每个案例有一个观测条件，它只影响上表字段的读数，不改变业务规则和标签。标准条件下的读数规律与
\texttt{train}、\texttt{cal\_fit}、\texttt{cal\_select}
相同；偏移条件下，这些读数的可靠性会发生变化，具体变化方式不公开，不同偏移案例之间也不完全相同。\texttt{dev}
中有 60
道带标签的偏移条件案例，使用了它们的方法不能声称没有用过偏移条件的标签。数据中不标注单个案例的条件。

\textbf{数据源与读取方式。} 返回有噪声字段的工具可以在请求中指定
\texttt{mode}。每个业务域有两个数据源：\texttt{default}、\texttt{refresh}
和 \texttt{remeasure} 读取主数据源，\texttt{independent}
读取第二数据源。名字只说明读取的是哪个数据源，不保证两者的误差统计独立。回执中的
\texttt{source\_id}
标识数据源，同一业务域的所有案例共用这两个标识；同一数据源的读数共享该数据源的误差特性，不能当作相互独立的证据。

{\small\def\LTcaptype{none} % do not increment counter
\begin{longtable}[]{@{}
  >{\raggedright\arraybackslash}p{(\linewidth - 6\tabcolsep) * \real{0.1842}}
  >{\raggedright\arraybackslash}p{(\linewidth - 6\tabcolsep) * \real{0.1316}}
  >{\raggedright\arraybackslash}p{(\linewidth - 6\tabcolsep) * \real{0.0789}}
  >{\raggedright\arraybackslash}p{(\linewidth - 6\tabcolsep) * \real{0.6053}}@{}}
\toprule\noalign{}
\begin{minipage}[b]{\linewidth}\raggedright
\texttt{mode}
\end{minipage} & \begin{minipage}[b]{\linewidth}\raggedright
数据源
\end{minipage} & \begin{minipage}[b]{\linewidth}\raggedright
费用
\end{minipage} & \begin{minipage}[b]{\linewidth}\raggedright
返回的记录
\end{minipage} \\
\midrule\noalign{}
\endhead
\bottomrule\noalign{}
\endlastfoot
\texttt{default} & 主数据源 & 0.01 &
可能是之前存下的旧记录；同一字段、同一阶段重复请求总是返回同一条记录 \\
\texttt{refresh} & 主数据源 & 0.01 & 不使用 \texttt{default}
存下的记录，但不保证是本次新采的记录。危险品分数和飞机的四个字段在每个数据源、每个阶段只有一条当前记录，重复请求返回同一条；客服指标重复请求可以得到新读数，其中第一次与
\texttt{default} 不用旧记录时的读数相同 \\
\texttt{remeasure} & 主数据源 & 0.01 & 目前与 \texttt{refresh}
完全相同，两者共用编号：第 n 次 \texttt{remeasure} 与第 n 次
\texttt{refresh} 返回同一条记录 \\
\texttt{independent} & 第二数据源 & 0.03 &
返回第二数据源在本阶段的当前记录。危险品分数和飞机的四个字段每个阶段只有一条；客服指标从第二次请求起每次返回新读数 \\
\end{longtable}
}

其他工具不受 \texttt{mode} 影响，任何方式都按 0.01 收费。

\textbf{回执字段。} \texttt{observation\_id}
是记录编号：两个读数的编号相同，说明是同一条记录，第二次读取没有带来新信息。\texttt{sampled\_at}
是记录的采样时间，\texttt{observed\_at}
是本次请求的业务时间；\texttt{sampled\_at} 早于
\texttt{observed\_at}，说明返回的不是本次采样的记录。客服的诊断和排障工具一次返回三个指标，每个指标的
\texttt{field}、\texttt{source\_id}、\texttt{sampled\_at} 和
\texttt{observation\_id} 列在回执的 \texttt{observations} 中。

\textbf{其他说明。}

\begin{itemize}
\tightlist
\item
  SOP、工具说明、输出格式和业务规则随代码提供；每个案例实际可用的工具和参数以服务返回为准。工具请求只能包含工具
  schema 中声明的参数；oracle 只接受服务规定的字段和业务阶段。
\item
  工具调用成功只表示这次调用完成，不表示模型的结论正确。客服的读数要按业务阶段解读，后一阶段的状态只有完成规定操作后才能读取。
\item
  \texttt{challenge\_data/} 另附 72
  个边界案例，包括恰好在阈值上、略低或略高于阈值、单项异常等情况；它们使用独立的数据版本，只用来检查规则实现，不计入测试成绩。
\item
  主数据是合成数据。去重会改变生成参数给出的比例；校准集和交互题按业务分支覆盖挑选，其中的比例不代表真实场景的发生率。危险品任务考察的是根据工具给出的分数分类，不考察阅读真实
  SDS 文档的能力；飞机的 \texttt{aircraft\_ready}
  只表示行政和维修记录上的放行结果，不是真实的适航结论。
\end{itemize}

\section{6. 评分细节}\label{ux8bc4ux5206ux7ec6ux8282}

\textbf{有效性。} \(p_i\) 要求初始状态为 \texttt{completed}，\(r_i\)
要求终局状态为 \texttt{completed}，\(q_i\) 要求两者都为
\texttt{completed}。状态为 \texttt{failed} 或 \texttt{missing}
时，对应的答案按错误计，风险按缺失计，缺失风险的 Brier 项记为
1；初始状态失败不影响已完成终局的评分。服务拒绝非有限值或超出 \([0,1]\)
的概率。

\textbf{控制器失效。} 同一题累计有超过 5 次请求（工具、oracle
或答案提交）被服务拒绝时，该题记为控制器失效：三项风险都按缺失计，决策损失按
5 计，转人工也不能抵消。

\textbf{决策损失。} 每道交互题的损失为已发生的工具和 oracle 费用
\(c_i\)，加上：

\begin{itemize}
\tightlist
\item
  转人工：0.5。终局状态为 \texttt{completed}、\texttt{failed} 或
  \texttt{missing} 都按 0.5 计；此时提交的 \(b_i\) 和 \(r_i\)
  仍按上面的有效性规则评估 \(r_i\)。
\item
  自动处理且答案正确：0。
\item
  自动处理且为普通错误：1。
\item
  记 5 的情况：自动处理且为严重错误；自动处理但终局状态不是
  \texttt{completed}，或答案缺失、不合法；没有提交决策；控制器失效。
\end{itemize}

严重错误指：低估危险品等级，包括把应为 \texttt{Unable\ to\ Decide}
的案例判为某个等级；把不该放行的飞机判为
\texttt{aircraft\_ready=true}；把没解决的客服问题判为
\texttt{RESOLVED}。转人工后人工的处理结果不作为模型是否出错的标签。

\textbf{辅助指标。} 下列指标只使用有效概率，同时报告有效提交的比例（NLL
除外，它在全部案例上计算）。

\begin{itemize}
\tightlist
\item
  NLL：二元交叉熵，有效概率截断到 \([10^{-6},1-10^{-6}]\)，缺失项记为
  \(-\log 10^{-6}\)。
\item
  ECE 和 aECE：分箱后预测风险与实际错误率之差的加权平均。ECE 用 10
  个等宽分箱；aECE 按样本数自适应分箱，概率相同的样本不拆到不同分箱。
\item
  AP：以出错为正类，衡量按风险从高到低排序时能否先找出错误；风险相同的样本作为一组处理，没有正类时不计算。
\item
  AURC：按风险从低到高逐个纳入案例，对各覆盖率下的错误率取平均；风险相同的案例按随机顺序取期望。
\item
  自动处理比例、自动处理错误率、转人工比例和平均查询费用，用于解释决策损失。
\end{itemize}

\textbf{置信区间。} 主指标附 95\% bootstrap
分位数区间：在每个业务域内按案例重抽样 2000
次。这个区间只反映这批案例的抽样波动，不代表分布变化后的表现。报告正式成绩时，要同时给出初始准确率、自动处理比例和业务损失，不能只凭某一项
Brier score 判断模型的业务能力。

\section{7. 模型与 baseline
细节}\label{ux6a21ux578bux4e0e-baseline-ux7ec6ux8282}

\textbf{模型。} M9 为 \texttt{Qwen/Qwen3.5-9B}，BF16 权重；M27 为
\texttt{Qwen/Qwen3.8-27B} 的 \texttt{UD-Q4\_K\_M} GGUF
量化版本。正式考试默认使用 M9，M27
用于考官另行配置的扩展评测，两者的成绩分开报告，量化方式也属于配置的一部分。文件哈希、tokenizer、推理设置和运行方式见
\texttt{model\_profiles/}。可以用所附部署代码自己运行
Qwen，也可以连接考官提供的同协议公共服务；Qwen
服务看不到完整数据或真实状态。

\textbf{参考 agent。} 参考轨迹关闭 thinking，使用 greedy
解码。对给出合法初始答案的案例，在生成该答案时的同一份历史上再采样四次（temperature
0.7、top-p 0.9、top-k
40），供所有风险估计方法共用；采样中提出的工具调用不会执行。参考 agent
在初始阶段最多调用模型 20 次，继续作答后累计最多 32 次。这是参考 agent
和 baseline
的设置，考生可以用别的方法继续作答。自己部署模型时留下的调用记录用于复现，不能代替考官对模型运行过程的独立核查。

\textbf{baseline。}
风险估计包括模型自报概率、采样不一致率、训练集错误率，以及基于公开特征的监督模型和概率校准。策略包括不查询、全部转人工、按阈值转人工、按固定字段查询、查询当前所有可查的字段，以及根据已有数据挑选字段查询的启发式方法。它们只使用公开输入、合法获得的观测、模型输出和训练标签。

\begin{itemize}
\tightlist
\item
  所附校准文件记录了生成训练轨迹时的模型、提示词、输入整理方式、采样设置、调用上限、风险估计代码、公共协议代码、预算和费用。其中任何一项变化后，旧的校准文件不能直接使用，需要重新采集轨迹或重新拟合。自己实现的在线方法可以随时更新参数，但要记录更新过程。
\item
  查询后的两个风险，分别用训练数据上的查询轨迹单独校准，并且必须和采集这些轨迹时的查询策略一致。只有在实际要查询时才检查是否有匹配的校准文件；某题没有查询时，直接沿用查询前的风险。不经校准的自报方法在查询前后都使用模型自报的概率，单独比较。
\item
  转人工阈值在 \texttt{cal\_select} 或 \texttt{dev}
  上按实际错误代价选取，手动指定的阈值会标为启发式。普通错误和严重错误的代价不同，只看出错概率不能确定最优决策。
\item
  程序会报告全部转人工、风险值几乎相同等异常情况。按已有数据挑选字段的
  baseline 是启发式方法，不保证查询的信息价值最优。
\end{itemize}

\textbf{比较方法时的注意事项。} 在同一个固定初始 agent
上比较风险估计方法时，使用考官提供的同一批案例和初始轨迹，并报告每种方法新获得的数据和花费。使用其他题目的信息或在线更新的方法，预测时可用的信息与固定历史不同，应记录题目的处理顺序和学习过程；另做一个只用固定历史的对照，有助于解释收益来自哪里。后续求解的模型或算法可以更换，但要区分哪些收益来自更好的求解器，哪些来自更好的风险估计。多个方法或模型共用同一份测试调用记录，按模型分列的费用不能看作各自独立获取信息的成本。

\section{8.
禁止事项与报告要求}\label{ux7981ux6b62ux4e8bux9879ux4e0eux62a5ux544aux8981ux6c42}

\begin{enumerate}
\def\labelenumi{\arabic{enumi}.}
\tightlist
\item
  不得越权获取测试答案、读取其他考生的记录，或伪造调用、回执和结果来源。自己用模型或算法算出的答案和概率是合法的。
\item
  无法给出合法答案或有效概率时，如实提交失败或缺失状态，不得编造。
\item
  公开报告只能使用训练标签和公布后的测试汇总结果，不能包含测试题的
  oracle 返回值、服务凭据或其他考生的材料；复现材料单独交给考官。
\end{enumerate}

\end{document}
